\documentclass[12pt,a4paper]{article}
\usepackage[a4paper,margin=18mm]{geometry}
\usepackage[T2A]{fontenc}
\usepackage[utf8]{inputenc}
\usepackage[russian]{babel}
\usepackage{amsmath,amssymb,mathtools}
\usepackage{array,booktabs,longtable}
\usepackage{xcolor}
\usepackage{hyperref}
\usepackage{tikz}
\hypersetup{colorlinks=true,linkcolor=blue!55!black,urlcolor=blue!55!black}
\setlength{\parindent}{0pt}
\setlength{\parskip}{0.55em}
\setlength{\emergencystretch}{2em}
\renewcommand{\arraystretch}{1.25}
\newcolumntype{L}[1]{>{\raggedright\arraybackslash}p{#1}}
\begin{document}
\begin{titlepage}
\centering
\vspace*{0.22\textheight}
{\LARGE\bfseries Дополнение к ревью домашней работы\par}
\vspace{2.2cm}
{\large Автор: Лёвин Артём Александрович\par}
\vspace{0.8cm}
{\large 15 сентября 2026 года\par}
\vfill
\begin{tikzpicture}
\draw[blue!35!black,line width=0.8pt](0,0)..controls(2,0.45)and(4,-0.45)..(6,0)..controls(8,0.45)and(10,-0.45)..(12,0);
\end{tikzpicture}
\vspace{0.8cm}
\end{titlepage}
\newpage
\section*{Сводная таблица ошибок}
{\small
\setlength{\tabcolsep}{3pt}
\begin{longtable}{@{}L{1.5cm}L{2.8cm}L{4.9cm}L{2.3cm}L{2.6cm}@{}}
\toprule
\textbf{№ задания}&\textbf{Тема}&\textbf{Суть ошибки}&\textbf{Ваш ответ}&\textbf{Правильный ответ}\\
\midrule
\endfirsthead
\toprule
\textbf{№ задания}&\textbf{Тема}&\textbf{Суть ошибки}&\textbf{Ваш ответ}&\textbf{Правильный ответ}\\
\midrule
\endhead
\hyperref[task:16]{16}&Вычитание десятичных дробей&Записано другое выражение с обыкновенными дробями; исходная разность не вычислена.&$\frac{21}{8}$ (к другому выражению)&$0{,}764$\\[0.4em]
\hyperref[task:17]{17}&Сложение и деление дробей&Записано другое деление дробей; выражение в скобках из условия не вычислено.&$\frac{33}{4}$ (к другому выражению)&$2$\\[0.4em]
\hyperref[task:18]{18}&Вычитание смешанных чисел&Вместо разности выполнено сокращение дроби $\frac{84}{126}$.&$\frac23$ (к другому выражению)&$1\frac1{14}$\\[0.4em]
\hyperref[task:19]{19}&Деление на произведение дробей&Вместо данного выражения выполнено разложение числа $180$ на простые множители.&не указан&$5$\\[0.4em]
\hyperref[task:20]{20}&Сокращение дроби&Вместо $\frac{360}{504}$ записано другое вычисление; обязательное разложение числителя и знаменателя не выполнено.&$\frac{11}{36}$ (к другому выражению)&$\frac57$\\
\bottomrule
\end{longtable}}
\newpage
\thispagestyle{empty}
\vspace*{\fill}
\begin{center}{\LARGE\bfseries Разбор ошибок}\end{center}
\vspace*{\fill}
% =========================================================
% Задание 16
% =========================================================
\newpage
\phantomsection
\label{task:16}
\section*{Задание 16}

\textbf{Текст задания.}
Вычислите:
\[
1{,}75-0{,}986.
\]

\textbf{Ваше решение.}
В строке, помеченной номером 16, записано вычисление с обыкновенными дробями,
которое заканчивается результатом
\[
\frac{21}{8}.
\]
Эта запись не соответствует выражению с десятичными дробями из условия.

\textbf{Ваш ответ.}
Ответ к заданному выражению не получен; в строке записан результат
\[
\frac{21}{8}
\]
для другого вычисления.

\textcolor{red}{\textbf{Ошибка:}} решение начато не с того выражения, которое дано в условии.

\textbf{Где именно возникает ошибка.}
Последнего корректного шага, относящегося к заданию 16, в записи нет.
Первый достоверно неверный шаг возникает сразу: вместо разности
\[
1{,}75-0{,}986
\]
записано другое выражение с обыкновенными дробями.
Из-за этого дальнейший результат не может быть ответом на данное задание.

\textbf{Почему это неверно.}
При проверке вычислительного задания необходимо сначала точно переписать исходное выражение.
Здесь требуется вычитание десятичных дробей. Нужно выровнять запятые и одинаковые разряды:
\[
1{,}75=1{,}750.
\]
После этого вычитаем тысячные из тысячных, сотые из сотых и десятые из десятых.

\textcolor{blue!60!black}{\textbf{Ключевая идея:}}
сначала сохраняем исходное выражение без замены на другой пример, затем выравниваем десятичные дроби по запятой.

\textbf{Исправление.}
\[
1{,}750-0{,}986=0{,}764.
\]

\textbf{Полное правильное решение.}
\[
1{,}75-0{,}986
=1{,}750-0{,}986
=0{,}764.
\]

\textbf{Проверка.}
\[
0{,}764+0{,}986=1{,}750=1{,}75.
\]

\textcolor{teal!60!black}{\textbf{Итог:}}
\[
\boxed{0{,}764}.
\]

\textbf{Задача на закрепление.}
Вычислите:
\[
2{,}4-1{,}387.
\]

% =========================================================
% Задание 17
% =========================================================
\newpage
\phantomsection
\label{task:17}
\section*{Задание 17}

\textbf{Текст задания.}
Вычислите:
\[
\left(\frac{3}{5}+\frac{7}{10}\right):\frac{13}{20}.
\]

\textbf{Ваше решение.}
Под номером 17 записано другое деление дробей; вычисление заканчивается результатом
\[
\frac{33}{4}.
\]
Исходная сумма $\frac35+\frac7{10}$ в этой строке не вычисляется.

\textbf{Ваш ответ.}
Для заданного выражения ответ не получен; записан результат $\frac{33}{4}$ для другого выражения.

\textcolor{red}{\textbf{Ошибка:}} вместо исходного выражения решается другой пример.

\textbf{Где именно возникает ошибка.}
Последнего корректного шага по заданию 17 нет: уже первая записанная операция не совпадает с условием.
Первый неверный шаг --- подмена выражения
\[
\left(\frac35+\frac7{10}\right):\frac{13}{20}
\]
другим делением дробей.

\textbf{Почему это неверно.}
Сначала нужно выполнить действие в скобках, то есть сложить $\frac35$ и $\frac7{10}$,
а затем делить полученную сумму на $\frac{13}{20}$.

\textcolor{blue!60!black}{\textbf{Ключевая идея:}}
сначала действие в скобках, затем деление на дробь как умножение на обратную дробь.

\textbf{Исправление.}
\[
\frac35=\frac6{10},
\qquad
\frac6{10}+\frac7{10}=\frac{13}{10}.
\]
Тогда
\[
\frac{13}{10}:\frac{13}{20}
=\frac{13}{10}\cdot\frac{20}{13}
=2.
\]

\textbf{Полное правильное решение.}
\[
\left(\frac35+\frac7{10}\right):\frac{13}{20}
=\left(\frac6{10}+\frac7{10}\right):\frac{13}{20}
=\frac{13}{10}:\frac{13}{20}
=\frac{13}{10}\cdot\frac{20}{13}
=2.
\]

\textbf{Проверка.}
\[
2\cdot\frac{13}{20}=\frac{26}{20}=\frac{13}{10},
\]
то есть возвращаемся к значению суммы в скобках.

\textcolor{teal!60!black}{\textbf{Итог:}}
\[
\boxed{2}.
\]

\textbf{Задача на закрепление.}
Вычислите:
\[
\left(\frac23+\frac56\right):\frac34.
\]

% =========================================================
% Задание 18
% =========================================================
\newpage
\phantomsection
\label{task:18}
\section*{Задание 18}

\textbf{Текст задания.}
Вычислите:
\[
2\frac37-1\frac5{14}.
\]

\textbf{Ваше решение.}
Под номером 18 записано сокращение другой дроби:
\[
\frac{84}{126}=\frac23.
\]
Это вычисление не связано с разностью смешанных чисел из условия.

\textbf{Ваш ответ.}
Для задания 18 ответ не получен; записан результат другого задания: $\frac23$.

\textcolor{red}{\textbf{Ошибка:}} вместо заданной разности выполнено сокращение другой дроби.

\textbf{Где именно возникает ошибка.}
Последнего корректного шага по исходной разности нет.
Первый неверный шаг появляется сразу: выражение
\[
2\frac37-1\frac5{14}
\]
заменено дробью $\frac{84}{126}$.

\textbf{Почему это неверно.}
Чтобы вычесть смешанные числа, нужно работать с их целыми и дробными частями
или перевести оба числа в неправильные дроби. Здесь удобно привести знаменатели к $14$.

\textcolor{blue!60!black}{\textbf{Ключевая идея:}}
переводим смешанные числа в неправильные дроби и приводим их к общему знаменателю.

\textbf{Исправление.}
\[
2\frac37=\frac{17}{7}=\frac{34}{14},
\qquad
1\frac5{14}=\frac{19}{14}.
\]
Поэтому
\[
\frac{34}{14}-\frac{19}{14}=\frac{15}{14}=1\frac1{14}.
\]

\textbf{Полное правильное решение.}
\[
2\frac37-1\frac5{14}
=\frac{17}{7}-\frac{19}{14}
=\frac{34}{14}-\frac{19}{14}
=\frac{15}{14}
=1\frac1{14}.
\]

\textbf{Проверка.}
\[
1\frac1{14}+1\frac5{14}=2\frac6{14}=2\frac37.
\]

\textcolor{teal!60!black}{\textbf{Итог:}}
\[
\boxed{1\frac1{14}}.
\]

\textbf{Задача на закрепление.}
Вычислите:
\[
3\frac25-1\frac7{10}.
\]

% =========================================================
% Задание 19
% =========================================================
\newpage
\phantomsection
\label{task:19}
\section*{Задание 19}

\textbf{Текст задания.}
Вычислите:
\[
3\frac13:\left(\frac56\cdot\frac45\right).
\]

\textbf{Ваше решение.}
Под номером 19 записано разложение другого числа на простые множители:
\[
180=2^2\cdot3^2\cdot5.
\]
Эта запись сама по себе верна, но она не относится к выражению задания 19.

\textbf{Ваш ответ.}
Не указан.

\textcolor{red}{\textbf{Ошибка:}} выполнено действие из другого задания; исходное выражение не вычислено.

\textbf{Где именно возникает ошибка.}
Последнего корректного шага по заданию 19 нет.
Первый неверный шаг --- переход от данного выражения к разложению числа $180$, которого в условии нет.

\textbf{Почему это неверно.}
Сначала нужно вычислить произведение в скобках, затем разделить смешанное число на полученную дробь.
Разложение числа $180$ на простые множители здесь не используется.

\textcolor{blue!60!black}{\textbf{Ключевая идея:}}
сначала выполняем умножение в скобках, затем переводим смешанное число в неправильную дробь
и деление заменяем умножением на обратную дробь.

\textbf{Исправление.}
\[
\frac56\cdot\frac45=\frac46=\frac23,
\qquad
3\frac13=\frac{10}{3}.
\]
Тогда
\[
\frac{10}{3}:\frac23
=\frac{10}{3}\cdot\frac32
=5.
\]

\textbf{Полное правильное решение.}
\[
3\frac13:\left(\frac56\cdot\frac45\right)
=\frac{10}{3}:\frac23
=\frac{10}{3}\cdot\frac32
=5.
\]

\textbf{Проверка.}
\[
5\cdot\frac23=\frac{10}{3}=3\frac13.
\]

\textcolor{teal!60!black}{\textbf{Итог:}}
\[
\boxed{5}.
\]

\textbf{Задача на закрепление.}
Вычислите:
\[
2\frac12:\left(\frac34\cdot\frac23\right).
\]

% =========================================================
% Задание 20
% =========================================================
\newpage
\phantomsection
\label{task:20}
\section*{Задание 20}

\textbf{Текст задания.}
Сократите дробь
\[
\frac{360}{504},
\]
предварительно разложив числитель и знаменатель на простые множители.

\textbf{Ваше решение.}
Под номером 20 записано другое вычисление с обыкновенными дробями,
которое заканчивается результатом
\[
\frac{11}{36}.
\]
Разложения чисел $360$ и $504$ в этой строке нет.

\textbf{Ваш ответ.}
Для дроби $\frac{360}{504}$ ответ не указан; записано $\frac{11}{36}$ как результат другого выражения.

\textcolor{red}{\textbf{Ошибка:}} выполнено другое вычисление и пропущено обязательное разложение числителя и знаменателя на простые множители.

\textbf{Где именно возникает ошибка.}
Последнего корректного шага по заданию 20 нет.
Первый неверный шаг --- замена дроби $\frac{360}{504}$ другим выражением.
Кроме того, условие требует сначала разложить оба числа на простые множители.

\textbf{Почему это неверно.}
Сокращение дроби означает деление числителя и знаменателя на один и тот же общий множитель.
Разложение на простые множители позволяет увидеть весь общий множитель сразу.

\textcolor{blue!60!black}{\textbf{Ключевая идея:}}
разложить числитель и знаменатель на простые множители, выделить одинаковые множители и сократить их.

\textbf{Исправление.}
\[
360=2^3\cdot3^2\cdot5,
\qquad
504=2^3\cdot3^2\cdot7.
\]
Тогда
\[
\frac{360}{504}
=\frac{2^3\cdot3^2\cdot5}{2^3\cdot3^2\cdot7}
=\frac57.
\]

\textbf{Полное правильное решение.}
\[
360=36\cdot10=(2^2\cdot3^2)(2\cdot5)=2^3\cdot3^2\cdot5,
\]
\[
504=8\cdot63=2^3\cdot3^2\cdot7.
\]
Следовательно,
\[
\frac{360}{504}
=\frac{2^3\cdot3^2\cdot5}{2^3\cdot3^2\cdot7}
=\frac57.
\]

\textbf{Проверка.}
Общий множитель равен
\[
2^3\cdot3^2=72.
\]
Действительно,
\[
360:72=5,
\qquad
504:72=7.
\]
Числа $5$ и $7$ взаимно простые.

\textcolor{teal!60!black}{\textbf{Итог:}}
\[
\boxed{\frac57}.
\]

\textbf{Задача на закрепление.}
Сократите дробь
\[
\frac{420}{588},
\]
предварительно разложив числитель и знаменатель на простые множители.


\vfill
\begin{center}
\small Лёвин Артём Александрович эксклюзивно для Ксении Васильченко
\end{center}
\end{document}
